%--------------------
% Packages
% -------------------
\documentclass[11pt,a4paper]{article}
\usepackage[utf8x]{inputenc}
\usepackage[T1]{fontenc}
%\usepackage{gentium}
\usepackage{mathptmx} % Use Times Font


\usepackage[pdftex]{graphicx} % Required for including pictures
\usepackage[swedish]{babel} % Swedish translations
\usepackage[pdftex,linkcolor=black,pdfborder={0 0 0}]{hyperref} % Format links for pdf
\usepackage{calc} % To reset the counter in the document after title page
\usepackage{enumitem} % Includes lists

\frenchspacing % No double spacing between sentences
\linespread{1.2} % Set linespace
\usepackage[a4paper, lmargin=0.1666\paperwidth, rmargin=0.1666\paperwidth, tmargin=0.1111\paperheight, bmargin=0.1111\paperheight]{geometry} %margins
%\usepackage{parskip}

\usepackage[all]{nowidow} % Tries to remove widows
\usepackage[protrusion=true,expansion=true]{microtype} % Improves typography, load after fontpackage is selected

\usepackage{lipsum} % Used for inserting dummy 'Lorem ipsum' text into the template


%-----------------------
% Set pdf information and add title, fill in the fields
%-----------------------
\hypersetup{ 	
pdfsubject = {},
pdftitle = {Exercise 2 - Revising For Flow},
pdfauthor = {}
}

%-----------------------
% Begin document
%-----------------------
\begin{document} %All text i dokumentet hamnar mellan dessa taggar, allt ovanför är formatering av dokumentet

\section{Exercise 2: Revising For Flow}

This exercise is intended to teach strategies for identifying problems with flow, transition and organization across writing.

\paragraph{} Your tasks for this exercise are the following:
\begin{itemize}
\item Read the Introduction section.
\item Re-order the paragraphs/ideas in the most logical order.You may need to move certain sentences around too. Separate the text into paragraphs or merge them into one if it helps to highlight logic or ideas of your narration.
\item Do the paragraphs connect and transition well?
\end{itemize}


Note: there is no need for additions or intense editing (grammar/spelling) - at least for this exercise! If you are running into any problems with Overleaf or have questions, please feel free to raise your hand and ask for help.

\subsection*{Introduction}

The purpose of this study was to examine whether the level of species diversity of non-woody wildflower plant communities adjacent to the path differed from that of the surrounding field, within the University of Toronto Mississauga’s Old Field. 

Factors that influence the level of diversity include natural factors (such as climate change and the spread of invasive species) and human disturbances (Krebs 1978). 

Biodiversity may be an abstract concept, but it can be seen in action regularly, through examples such as wildflower plants that propagate the survival of various species by acting as a nectar source for insects, a food source for animals and facilitating the pollination of other plant species (Korpela et al. 2013). 

In fact, human actions, such the loss and fragmentation of habitats, overfishing, excessive pesticide use, trampling and irrigation, have led to a negative impact on biodiversity (Krebs 1978). When plant community structures are altered through such actions, these disturbances disrupt present biotic interactions and result in a series of significantly negative ecological impacts (Bond 1994). 

An ecosystem’s level of biodiversity can be monitored through ecological characteristics such as species richness, relative abundance and diversity indices (Krebs 1978). 

Queiroz et al. (2014) investigated whether the species composition of plant communities adjacent to hiking trails differed significantly from that of the surrounding vegetation in the Azores. While the trail had a slight impact on species composition, the authors concluded that there were no significant changes in community composition between the two sites (Queiroz et al. 2014). On the other hand, Roovers et al. (2004) evaluated plant species variation (by recording species richness) across both path and undisturbed areas of sampled transects for four different sites. The authors concluded that species diversity was linked to the species structure of the area examined – specifically that species diversity was greatest at the trail centre (Roovers et al. 2004). 

Given previous literature, it is expected that there will be a difference in the level of species diversity between the two plant communities, specifically that the level of species diversity will be higher at the field site, as this field site experiences lesser human disturbance.

Similarly, Behera (2010) investigated plant diversity across 49 sample plots in the Eastern Himalaya landscape. Results indicated that forests with greater fragmentation had fewer plant species, leading the author to conclude that species diversity was inversely proportional to habitat fragmentation (Behera 2010).

\end{document}
